% Paper B: discussion and conclusion.
\section{Discussion}\label{sec:discussion}
\subsection{What is proved, what is measured}
\Cref{tab:status} separates the three kinds of statement in this paper. The theorems hold for the stated hypotheses; the
algorithm uses them as design rules, not as guarantees about its output; and the numerical results are evidence on
synthetic data with white noise.

% Shared: the kinds of claim (proved / design rule / measured).
\begin{table}[tbp]
\centering\small
\caption{The kinds of claim made in the paper.}\label{tab:status}
\begin{tabularx}{\linewidth}{@{}>{\raggedright\arraybackslash}p{2.2cm}X@{}}
\toprule
Kind & Statements\\
\midrule
proved & groups are factors with overlap $\deltaP$ (\cref{thm:factor}); terms are not groups (\cref{prop:terms,prop:discont});
identification criterion and its limits on finite data (\cref{thm:ident,prop:unique,prop:certify}); stability via
$\gN$ (\cref{thm:hm}); Bayesian error bars via the angle (\cref{prop:bayes-angle}); counting condition (\cref{prop:counting})\\
design rule & labels instead of ``identified''; distributions of $\rho$ instead of values; decisions on the raw observations;
growth of the class instead of coefficient penalties; minimum-norm split of the overlap\\
measured & $0/20$ confident wrong answers on random equations not used in development; $60/60$ factor sets on grid
designs with independent fields and noise; $30/30$ on the development benchmark; coefficient coverage $0.94$ at
nominal $0.90$\\
\bottomrule
\end{tabularx}
\end{table}


\subsection{Limits of inference}\label{sec:limits}
% Shared module: which inference problems are solvable and which are not.
The statistical layer inherits limits from the analysis. Some are engineering problems with a known
remedy; others are theorems and no amount of data or modelling removes them.

\paragraph{Problems with a remedy}
\begin{itemize}
  \item \emph{Choosing $k$.} Rounding $\keff$ is unstable; the posterior $P(k)$ replaces it.
  \item \emph{Correlated noise after interpolation.} Residuals of a Gaussian-process interpolant are
  correlated, so white-noise likelihoods miscalibrate the intervals. The remedy is generalized least squares
  per Fourier bin, with the bin covariance estimated from the posterior draws; the angle should then be
  measured in the inner product of the inverse posterior covariance, which the method approximates bin by bin.
  \item \emph{Matching factors across samples.} A nearest-vector match without a gate merged transport and
  heat; Dirichlet-process clusters with a co-occurrence merge rule fixed it.
  \item \emph{Real fields.} A complex-valued component model fitted to a real field undercovered; a real
  parametrization by conjugate pairs restored nominal coverage.
\end{itemize}

\paragraph{Problems without a remedy}
\begin{itemize}
  \item \emph{Error floor.} In the Gaussian model~\eqref{eq:linmodel} with a flat prior, every unbiased estimator of the
  component coefficients has variance at least $\noise^2/\gN$ along the worst direction (the Cram\'er--Rao bound: $\noise^2$ times the
  inverse of the smallest nonzero eigenvalue of $\sumop^H\sumop$, \cref{prop:bayes-angle}), and the minimum-norm split attains
  the corresponding amplification $\gN^{-1/2}$ (\cref{thm:hm}):
  closely aligned groups are intrinsically hard to separate.
  \item \emph{No data-free threshold.} Whether a numerical symbol ``factors'' depends on the noise level;
  there is no threshold independent of the data, which is why the decision is by evidence and not by a
  tolerance on coefficients.
  \item \emph{Repeated factors.} Finite noisy data cannot distinguish $P_i^2$ from two nearby factors;
  the label is \UND.
  \item \emph{Loss of a derivative.} Components may be one derivative rougher than $u$; the prior on the
  components must be one order smoother than the prior on the field.
  \item \emph{No certificate.} By \cref{sec:certify} no finite data certify an equation; outputs are
  compatibility statements with uncertainty.
\end{itemize}


\subsection{Beyond the linear ring}\label{sec:beyond}\label{sec:nonlinear}
% Shared module: what changes for nonlinear equations (short, clearly an outlook).
The results above rest on one structure: candidate equations form the commutative ring
$\ring=\C[\tau,\xi]$, solutions form vector spaces, and the exact sequence behind \cref{thm:factor} (a Chinese
remainder theorem up to the finite overlap $\deltaP$) turns a
factorization of the symbol into a splitting of the solution. For a nonlinear equation
$u_t=\sum_n\mathcal A_n(u)$, with each $\mathcal A_n$ a polynomial in the jet $(u,u_x,u_{xx},\dots)$, this structure is lost
in three ways.
\begin{enumerate}[label=(\roman*)]
  \item \emph{No ring of candidates.} Products of nonlinear operators are not operators of the same kind
  in any useful sense, and there is no unique factorization. One fixes instead a finite
  \emph{dictionary} $\mathcal D$ of jet monomials of bounded degree and order, with span $\mathcal V$, and asks
  which relations $F(u,u_x,\dots)=0$ in $\mathcal V$ the data satisfy. In practice each monomial (for instance
  $u_t$, $u_x$, $u_{xx}$, $uu_x$, $u^2$) is evaluated at the sample points, giving a matrix $\Theta$ with one column per
  monomial; a relation $F=\sum_mc_m\,m(u)$ is a vector $c$ with $\Theta c\approx0$, so the candidate equations are the
  near-null space of $\Theta$.
  \item \emph{No superposition.} Solutions do not form a vector space, so ``the component of $u$ along a
  factor'' has no meaning; the groups of \cref{thm:factor} do not exist.
  \item \emph{Identification is relative.} Identification can only hold relative to $\mathcal V$ and modulo the
  differential consequences of $F$ (its derivatives $\partial_xF,\partial_x^2F,\dots$ are also relations in
  any rich dictionary). The precise statement is an exact sequence relating the relations of $u$ in $\mathcal V$
  to those generated by $F$; identification holds exactly when a jet-level obstruction vanishes.
\end{enumerate}

What survives is the interaction term. Definition~\eqref{eq:eint-def} still makes sense with the
nonlinear single-term flows, and its short-time expansion is
\begin{equation}\label{eq:jet}
  \Eint(t)=\frac{t^2}{2}\sum_{n\ne m}\mathcal A_n'[u_0]\big(\mathcal A_m(u_0)\big)+O(t^3),
\end{equation}
where $\mathcal A_n'[u_0]$ is the Fr\'echet derivative of $\mathcal A_n$ at $u_0$. For Burgers' equation
$u_t=-uu_x+\nu u_{xx}$ the two terms are $\mathcal A_1(u)=-uu_x$ and $\mathcal A_2(u)=\nu u_{xx}$, with Fr\'echet
derivatives $\mathcal A_1'[u](v)=-(vu_x+uv_x)$ and $\mathcal A_2'[u](v)=\nu v_{xx}$. Then
\[
  \mathcal A_1'[u_0]\big(\mathcal A_2(u_0)\big)=-\nu\,(u_0''u_0'+u_0u_0'''),\qquad
  \mathcal A_2'[u_0]\big(\mathcal A_1(u_0)\big)=-\nu\,(3u_0'u_0''+u_0u_0'''),
\]
and half their sum gives
\[
  \Eint(t)=-\nu\,\big(u_0u_0'''+2u_0'u_0''\big)\,t^2+O(t^3).
\]
The expansion holds only while the solution is smooth. The flow of the single term $-uu_x$ is the inviscid
Burgers equation, whose characteristics cross: for $u_0=\sin x$ they are $x=\xi+t\sin\xi$, and
$\partial x/\partial\xi=1+t\cos\xi$ first vanishes at $\xi=\pi$, $t^*=1$. After this breaking time the single-term flow
has a jump, $\Eint$ has a jump in $x$ whose height does not vanish as $\nu\to0$, and the interaction can only be measured in
$L^1$, where for a single shock it is small, of order $\nu$ up to logarithmic factors (a rate we do not prove here).
Whether such data still identify the equation relative to a dictionary depends on how many shocks they contain
compared with the degree of the dictionary; we do not pursue this here.

\paragraph{Variable coefficients}
If the coefficients depend on $(t,x)$, the operators form the Weyl algebra $\C[t,x]\langle\partial_t,\partial_x\rangle$, which
is not commutative ($\partial_x\,x-x\,\partial_x=1$). Annihilators become left ideals, factorizations are no longer unique,
and \cref{thm:factor} has no general analogue. Worse, if the coefficients may be arbitrary smooth functions, one solution
identifies nothing: any $u$ without zeros satisfies the first-order equation $\big(\partial_x-u_x/u\big)u=0$, whose
coefficient is built from $u$ itself. Identification becomes possible only inside a fixed, finite-dimensional class of
coefficients. For the Schr\"odinger equation $iu_t=-u_{xx}+V(x)u$, for instance, the potential is read off directly as
$V(x)=(iu_t+u_{xx})/u$ wherever $u\neq0$. In the spectral picture of \cref{sec:alg-pencil}, slowly varying coefficients
would make the rates depend on position, $\lambda(\kappa;x)$ (\cref{rem:varcoef}).

\paragraph{A possible route} The grouping idea can be carried over by replacing ``modes on the same
irreducible component of $\Var(P)$'' with ``jet samples satisfying the same relation in $\mathcal V$''. Locally
in time, \eqref{eq:jet} shows that a nonlinear flow is a linear one (its Fr\'echet derivative) plus a
quadratic correction, so one can group linearized modes on short windows and test, through the size
of $\Eint$, whether the groups persist. We do not pursue this here: without a ring there is no
analogue of $\deltaP$, of the projectors~\eqref{eq:lagrange} or of $\gN$, and these are what make the
linear algorithm verifiable.

\begin{table}[tbp]
\centering\small
\caption{What carries over beyond scalar linear equations with constant coefficients. ``Groups'' means a splitting of
a solution into pieces attached to pieces of the equation.}\label{tab:beyond}
\begin{tabularx}{\linewidth}{@{}>{\raggedright\arraybackslash}p{2.3cm}>{\raggedright\arraybackslash}X>{\raggedright\arraybackslash}X>{\raggedright\arraybackslash}X@{}}
\toprule
Setting & Candidate equations & Groups & What survives\\
\midrule
linear, constant coefficients & commutative ring $\C[\tau,\xi]$; ideals; unique factorization & factors of the symbol, overlap $\deltaP$ (\cref{thm:factor}) & everything in this paper\\
variable coefficients & non-commutative Weyl algebra; left ideals; factorization not unique & no general splitting theorem & identification inside a finite-dimensional class of coefficients; local rates $\lambda(\kappa;x)$ (\cref{rem:varcoef})\\
systems (e.g.\ Dirac) & matrices over the ring; modules instead of ideals & sums of kernels need extra conditions and can fail on non-convex domains & identification criteria for square non-singular systems\\
nonlinear & finite dictionary of jet monomials; differential ideal & none (no superposition) & identification relative to the dictionary; the interaction term and its expansion~\eqref{eq:jet}\\
\bottomrule
\end{tabularx}
\end{table}


\paragraph{Extensions}
The grouping principle is a statement about the ring of constant-coefficient operators, and its reach is set by that
ring. For variable coefficients the eigenvalues read by the pencil become functions of position (\cref{rem:varcoef}) and
each group becomes a family of curves; the operators then form a non-commutative algebra in which factorization is not
unique, and a splitting theorem would have to be proved for a restricted class of coefficients. For systems, such as the
Dirac equation, the ideal becomes a module and splitting requires topological conditions beyond simple connectivity. For
nonlinear equations groups can only be defined relative to a dictionary, as discussed above. In each direction the
first task is the one carried out here for the linear case: decide which pieces of the equation a solution carries, how
uniquely and how stably, and only then design the learning procedure around them.

\section{Conclusion}\label{sec:conclusion}
Groups are factors. This one statement, made precise by \cref{thm:factor} and its companions, determines what an
equation-learning algorithm should look for (modes on the zero sets of factors), how it should combine evidence
(clustering factors across posterior samples), how it should measure uncertainty (through the Friedrichs angle), and
what it may claim (compatibility, never certification). The resulting method gives no confident wrong answer on twenty random equations not used in its
development, against thirteen or more for each of the three library-regression baselines run on that set, gives the right factor
set on all sixty grid inputs with independent fields and noise, and recovers every factor set of its
thirty-input development benchmark, including the products of factors on which regression over a library of terms fails.
